\chapter{Experiments}
\label{ch:experiments}

This chapter describes every experiment carried out, in the order in which the questions
arose. It is deliberately complete rather than selective: configurations that were
rejected are reported alongside those that were adopted, because several of the rejections
are the most informative results in the work. Chapter~\ref{ch:results} then reports what
the experiments produced.

Section~\ref{sec:protocol} fixes the protocol that every experiment obeys, since the
protocol is what makes the individual results comparable. Sections~\ref{sec:exp-arch}
to~\ref{sec:exp-tuning} then work through the five groups of experiments:
architecture selection, covariate selection, forecast reconciliation, forecast horizon,
and hyperparameter tuning.

%---------------------------------------------------------------------------------------
\section{Experimental protocol}
\label{sec:protocol}

\subsection{Why a single validation block was not enough}

The first version of this work selected covariates on one contiguous validation block.
That turned out to invert the answer. Eleven candidate feature sets were ranked on the
validation block and then re-ranked on the test period, and \emph{the two orderings were
completely reversed}: the most valuable family, per-station weather at the generating
plants, ranked eighth of eleven on validation and first on test.

The cause was seasonal. The validation block spanned summer, and almost all of the value
of weather information sits in winter, when wind and inflow drive dispatch. A single
block therefore measures the season it happens to cover, not the covariate.

Every experiment reported below consequently uses four-fold cross-season validation.

\subsection{Four-fold cross-season validation}

Each fold trains on an expanding window and evaluates on the three months that follow it.
The four evaluation windows are spaced one quarter apart so that between them they cover
all four seasons:

\begin{table}[htp]
\centering
\begin{tabular}{clll}
\toprule
{\bf Fold} & {\bf Train window ends} & {\bf Evaluation window} & {\bf Training origins} \\
\midrule
1 & 2025-04-15 & 2025-04-15 to 2025-07-15 & 3,114 \\
2 & 2025-07-15 & 2025-07-15 to 2025-10-15 & 3,296 \\
3 & 2025-10-15 & 2025-10-15 to 2026-01-15 & 3,480 \\
4 & 2026-01-15 & 2026-01-15 to 2026-04-15 & 3,664 \\
\bottomrule
\end{tabular}
\caption[Cross-validation folds]{\textbf{The four cross-validation folds.} All four
evaluation windows end on or before 2026-04-15, the first day of the test period, so no
fold can see test data. Training windows expand rather than slide, which is why the origin
count grows.}
\label{tab:folds}
\end{table}

Training origins are taken every twelve hours within the training window; evaluation
origins are taken once per day, and only where all 48 forecast steps fall inside the
evaluation window. Each origin expands into $48 \times 19 = 912$ rows, one per forecast
step per node.

\subsection{The decision rule: sign consistency, not the mean}
\label{sec:signrule}

A configuration is accepted only if it improves on the reference in \emph{all four folds}.
The four-fold mean is reported but is not the criterion.

The reason is that the per-fold differences are small relative to their own scatter.
Repeated measurements established a noise floor of approximately
$\pm 0.6$~\pp: below that, a single fold's number carries no interpretable information.
Several configurations produced four-fold means inside that floor while disagreeing in
sign between folds, and those are exactly the cases where a mean would mislead.

Sign consistency across four independent seasons is a sign test: under a null hypothesis
of pure noise, four same-signed differences occur with probability $1/16$. Its power does
not depend on the size of the effect, which is what makes it usable at this effect size.
Where a configuration passes the sign test but its mean falls inside the noise floor, that
is stated explicitly and the conclusion is confirmed independently on the test period.

\subsection{Two aggregations that must never be mixed}
\label{sec:aggregation}

Every percentage in this work is a relative change in mean absolute error against a stated
reference. Two different aggregations appear, and they answer different questions:

\begin{table}[htp]
\centering
\begin{tabular}{p{4.2cm}p{6.3cm}r}
\toprule
{\bf Where it is used} & {\bf How errors are aggregated} & {\bf Example} \\
\midrule
Cross-validation tables & Per node first, then the 19 nodes weighted equally & --- \\
Test-period tables & All errors pooled into one MAE; large nodes dominate & 9,240 \\
\bottomrule
\end{tabular}
\caption[The two aggregations]{\textbf{The two aggregations.} For the delivered model the
pooled MAE is 9{,}240~kWh per half hour while the median of the 19 per-node MAEs is
4{,}879 --- a factor of 1.9. Three large aggregate nodes account for the gap. Quoting one
number against the other would look like a large effect where there is none.}
\label{tab:aggregation}
\end{table}

\subsection{Three reference points}

Results are reported against three references, chosen so that each isolates a different
claim:

\begin{enumerate}
\item \textbf{\naiveA{}} --- the value observed 48 half-hours (24 hours) earlier. This is
the operational fallback and the denominator of MASE.
\item \textbf{The same model with no covariates} --- isolates the contribution of external
information from the contribution of the model class.
\item \textbf{That round's own baseline configuration} --- isolates the single change under
test, with every other choice held fixed.
\end{enumerate}

A fourth reference, \naiveB{} (the value one week earlier), is reported in the final
tables because it is the harder benchmark for a weekly-seasonal series.

\subsection{Uncertainty: paired bootstrap by forecast origin}
\label{sec:bootstrap}

The test period supplies 107 forecast origins. Each origin is an approximately independent
observation of forecast skill, so a single point estimate understates what is known and
overstates what is settled.

For every comparison, the errors within an origin are first averaged, giving 107 paired
observations; those pairs are then resampled 2,000 times. Pairing matters: two methods
compared on the same origin face the same day, so pairing removes the shared
``was that day easy or hard to forecast'' component, which is the dominant source of
variation. An interval that excludes zero is reported as a reliable difference; one that
straddles zero is reported as evidence insufficient, whatever the point estimate.

\subsection{Leakage discipline}

Covariates are named by the information they carry rather than by their source, so that a
leak is visible in the feature name:

\begin{itemize}
\item \verb+_at_origin+ --- the value known at the moment the forecast is issued;
\item \verb+_at_target+ --- the value at the time being forecast, admissible only for
quantities that are genuinely forecast and published ahead of time, such as numerical
weather prediction fields;
\item \verb+_change+ --- the difference between the two above;
\item \verb+_vs_yesterday+ --- the change against the same clock time one day earlier.
\end{itemize}

Families that are only ever observed, never forecast, are marked \verb+ORIGIN_ONLY+ and
may not appear with an \verb+_at_target+ suffix. This is enforced in the feature builder
rather than left to discipline.

A second rule follows from the model class. Gradient-boosted trees cannot extrapolate
beyond the range of a feature seen in training. Level quantities are therefore dangerous
whenever the future range can exceed the historical one --- a newly commissioned plant is
the obvious case --- whereas ratios and changes stay inside their historical range. Where
a level is needed it enters as a ratio to a long-run reference.

%---------------------------------------------------------------------------------------
\section{Architecture selection: nine candidates and the expert system}
\label{sec:exp-arch}

\subsection{Provenance of the code}

The expert system was written by a fellow student, Guoping, and made available through
GitHub. It was not modified for the experiments in this section: the code was run as
supplied, on the same node hierarchy, the same test period, and the same 107 forecast
origins as the pipeline of Section~\ref{sec:exp-cov}, so that the two systems are directly
comparable. Every number reported for it in this dissertation was produced by running that
code; the design of the system is his, and the analysis and discussion of its behaviour
are mine.

\subsection{What the system does}

The system trains nine candidate forecasters and then, for each of the 19 nodes, selects
the single candidate that performed best on a validation period. A tenth candidate,
\naiveA{}, is always available without training. The per-node choices are then passed
through a learned bottom-up network (referred to below as BUN) which restores hierarchical
coherence.

\begin{table}[htp]
\centering
\begin{tabular}{ll}
\toprule
{\bf Candidate} & {\bf Family} \\
\midrule
THP-Former & transformer, hierarchy-aware \\
PatchTST & transformer, patched input \\
TSMixer & all-MLP mixer \\
TiDE & dense encoder--decoder \\
NLinear & linear with normalisation \\
TCN & dilated temporal convolution \\
RNN, LSTM, GRU & recurrent \\
\midrule
\naiveA{} & no training; always eligible \\
\bottomrule
\end{tabular}
\caption[Candidate architectures]{\textbf{The nine trained candidate architectures},
spanning the transformer, MLP, convolutional and recurrent families, plus the untrained
seasonal-naive fallback.}
\label{tab:candidates}
\end{table}

\subsection{Experiments run on the expert system}

\begin{table}[htp]
\centering
\begin{tabular}{llp{6.1cm}}
\toprule
{\bf Run} & {\bf Window} & {\bf Question} \\
\midrule
\verb+expert_system+ & from 2023 & baseline behaviour \\
\verb+expert_system_2021+ & from 2021 & does a longer history help? \\
\verb+expert_system_2021_cov+ & from 2021 & do the 24 external covariates help? \\
\verb+diag_A_pocp+ & from 2021 & one clean covariate: plant availability \\
\verb+diag_B_ecm+ & from 2021 & one clean covariate: ECMWF weather \\
\bottomrule
\end{tabular}
\caption[Expert-system runs]{\textbf{The five expert-system runs.} The last two are
diagnostic: each supplies a single covariate column that is complete for all 19 nodes and
needs no imputation, in order to test whether the loss of accuracy under covariates was
caused by missing-value imputation.}
\label{tab:es-runs}
\end{table}

The 2021 run is the one used for all comparisons in Chapter~\ref{ch:results}, because it
matches the data window of the gradient-boosting pipeline. The diagnostic runs
\verb+diag_A_pocp+ and \verb+diag_B_ecm+ were added after the covariate runs produced a
result that contradicted the covariate findings of Section~\ref{sec:exp-cov}; their purpose
is described in Section~\ref{sec:exp-cov-es}.

%---------------------------------------------------------------------------------------
\section{Covariate selection: twenty-two cross-validation rounds}
\label{sec:exp-cov}

This is the largest group of experiments. Each round tests one question with every other
choice held fixed, and each round's baseline is re-run inside the same execution as its
variants. That last point is not a formality: an earlier pair of rounds that shared a
configuration name but were executed separately disagreed by 0.1 to 1.0~\pp{} on the
aggregate nodes, which is the same order as the effects being measured. Comparisons are
therefore only ever made within a single run.

\subsection{Constraints imposed by the supervisors}

Two constraints were given at the outset and are respected throughout:

\begin{enumerate}
\item Certain variables were specified by the supervisors as required, and are
\emph{locked}: they are never candidates for elimination, whatever a round reports about
them.
\item Different generator types use different variables. Results must therefore be read
per generator category, not as a median over all 19 nodes --- a variable that matters to
hydro and is irrelevant to geothermal will look like noise in a 19-node median.
\end{enumerate}

The second constraint is implemented as per-type gating. Each feature family declares the
node types it is allowed to reach; where a family does not apply to a node, the value is
masked as missing rather than set to zero, so that the model can distinguish
``not applicable'' from ``zero''. Two gating modes were tested: \verb+strict+, where a
family reaches only its own generator type, and \verb+dispatch+, where wind information is
also given to hydro because the two compete in the same dispatch stack.

\subsection{Inventory of rounds}

\begin{longtable}{clp{7.0cm}}
\caption[Cross-validation rounds]{\textbf{The twenty-two cross-validation rounds.}
Rounds 1--3 ran on the 2023 data window; rounds 4 onward ran on the extended 2021 window,
which required the no-covariate baseline to be re-run because the baseline also benefits
from the longer history.}
\label{tab:rounds} \\
\toprule
{\bf Round} & {\bf Configurations} & {\bf Question} \\
\midrule
\endfirsthead
\toprule
{\bf Round} & {\bf Configurations} & {\bf Question} \\
\midrule
\endhead
\bottomrule
\endfoot
1 & 7 & Which weather source, and does hydrology help at all? \\
2 & 4 & Company-level hydro, storage, lake level, and weather combined with it \\
3 & 5 & Seasonal terms, a one-year principal component, seven-day changes \\
4 & 6 & Repeat of the survivors on the longer 2021 history \\
5 & 3 & Full set, and per-node gating for the first time \\
6 & 2 & Gating in \verb+dispatch+ mode rather than \verb+strict+ \\
7 & 2 & Forecast horizon other than 48 steps (see Section~\ref{sec:exp-horizon}) \\
8 & 2 & Availability versus the raw POCP channel \\
9 & 4 & Supervisor-specified families, one at a time \\
10 & 1 & Those families combined \\
12 & 7 & First backward-elimination pass: remove one family at a time \\
13 & 4 & Second backward-elimination pass on the survivors \\
14 & 3 & Hydrology v4: which of its three groupings \\
15 & 4 & The two remaining supervisor-specified columns \\
16 & 4 & Aggregate-node weather split into wind and non-wind \\
17 & 3 & Hierarchy features (parent and sibling) ablated \\
18 & 5 & Full aggregate-node ladder, superseding rounds 16 and 18's predecessors \\
19 & 6 & What each stage of the selection process is worth \\
20 & 2 & What the manual per-type assignment is worth \\
21 & 4 & Five additional weather variables, added rather than removed \\
22 & 2 & South Island snowpack \\
\end{longtable}

Three features of this table are worth drawing out.

\textbf{Rounds 12 and 13 are backward elimination.} Rather than adding families and
keeping whatever improves, which rewards whichever family happens to be tested first, each
family is removed from the full set in turn and the damage measured. A family earns its
place if removing it hurts in all four folds.

\textbf{Round 18 replaced rounds 16 and its predecessor.} The two earlier rounds had been
executed separately, and their shared baseline configuration disagreed between runs. Round
18 puts the baseline and all four variants into a single execution so that the ladder is
internally consistent. The superseded rounds are retained in the code and in
Appendix~\ref{ch:appendix} rather than deleted.

\textbf{Round 17 exposed a discrepancy between the measured and the delivered model.} The
hierarchy features --- a node's parent and sibling values --- had never been enabled in any
cross-validation round, because the default in the round runner was off, whereas the
delivered model has them on. Rounds 1--16 therefore measured an 84-feature model while the
delivered model has 89 features. Round 17 was added to ablate them. This discrepancy is
recorded rather than papered over, and the hyperparameter search of
Section~\ref{sec:exp-tuning} was deliberately run with the features enabled, so that it
tunes the model that is actually delivered.

\subsection{Rounds 21 and 22: two additions that were rejected}
\label{sec:exp-negative}

Two rounds tested enlarging the candidate pool rather than pruning it, and both were
rejected. They are reported because the asymmetry matters: when the question is
``should this be added?'', insufficient evidence means \emph{do not add}, whereas when the
question is ``should this be removed?'', insufficient evidence means \emph{keep}.

Round 21 added five further weather variables --- cloud cover, radiation, dew point,
evapotranspiration and soil moisture --- in three combinations. Round 22 added a South
Island snowpack estimate. The snowpack experiment was prompted by a specific observation:
the existing snow family covered only the Taupo and Waikato catchments, which have had
essentially no snow since 2021, so that family had been acting as a constant plus a trend
rather than as snow information. The South Island catchments do accumulate snow, so the
family was rebuilt there using a degree-day melt model.

\subsection{Why the covariate experiments were also run on the expert system}
\label{sec:exp-cov-es}

Feeding the same 24 covariate columns to the expert system made it worse, which
contradicts Section~\ref{sec:exp-cov}. Two explanations were available: the covariates
themselves, or the path by which they enter that particular system.

The diagnostic runs in Table~\ref{tab:es-runs} discriminate between them. Each supplies a
single covariate column that is complete across all 19 nodes and therefore needs no
imputation at all. If imputation of missing values were the cause, these runs should be
unharmed. This is reported in Chapter~\ref{ch:results}.

%---------------------------------------------------------------------------------------
\section{Reconciliation}
\label{sec:exp-recon}

\subsection{The four reconcilers}

All four reconcilers used here are linear projections of the base forecast vector and can
be written in one form,
\begin{eqnarray}\label{eqn:recon}
\hat{y}_{\mathrm{rec}} = S\left(S^{\top} W^{-1} S\right)^{-1} S^{\top} W^{-1}\, \hat{y},
\end{eqnarray}
where $S$ is the $19 \times 15$ summing matrix and the choice of $W$ distinguishes the
methods: the identity for OLS, the diagonal of residual variances for WLS, and the full
residual covariance for MinT. Bottom-up is the degenerate case in which only the 15 leaf
forecasts are used and the aggregates are recomputed from them.

For MinT the covariance is estimated with Schäfer--Strimmer shrinkage towards its
diagonal, with the shrinkage intensity determined analytically from the residuals rather
than tuned. All reconciliation weights are estimated on validation-period residuals; the
test period is never used to fit them.

\subsection{Experiments on the reconciliation layer}

\begin{table}[htp]
\centering
\begin{tabular}{p{4.6cm}p{8.4cm}}
\toprule
{\bf Experiment} & {\bf Question} \\
\midrule
Four reconcilers, both systems & Which projection, and is the ordering the same for a
gradient-boosting base and a neural base? \\
Re-run after every base change & Does the ordering survive a change of base model, given
that the residual covariance is re-estimated? \\
Learned BUN versus classical MinT & Is a learned bottom-up network better than a
closed-form projection? \\
Horizon-grouped weights & Should $W$ be estimated separately for groups of forecast steps,
given that the error correlation structure changes with lead time? \\
Negative-forecast audit & How often does each projection produce negative generation, which
is physically impossible? \\
\bottomrule
\end{tabular}
\caption[Reconciliation experiments]{\textbf{Experiments on the reconciliation layer.}}
\label{tab:recon-exp}
\end{table}

The horizon-grouped experiment deserves its protocol spelled out, because selecting the
grouping granularity is itself a selection problem. A single $W$ currently covers all 48
forecast steps, pooling $105 \times 48 = 5{,}040$ residual vectors. But at one step ahead
the errors are dominated by persistence and are strongly aligned across nodes, whereas at
48 steps they are dominated by weather and dispatch uncertainty with a different
correlation structure. Five granularities were tested, from one pooled $W$ down to one per
forecast step.

The granularity was selected using the validation period alone, split into two halves
chronologically: $W$ estimated on the first half and scored on the second, then the
reverse. The test period was used only once, afterwards. As reported in
Chapter~\ref{ch:results}, this selection design has a flaw that the results exposed, and
that flaw is reported rather than hidden.

%---------------------------------------------------------------------------------------
\section{Forecast horizon}
\label{sec:exp-horizon}

The supervisors asked whether the pipeline extends from ``the next day'' to ``the next two
weeks''. The horizon is a configuration parameter, and round 7 re-ran the no-covariate
baseline and the best covariate configuration at horizons other than 48 steps.

Because each origin expands into $\mathrm{horizon} \times 19$ rows, a fourteen-day horizon
produces fourteen times the rows per origin. The origin stride was widened in proportion
so that the total training set stays the same order of magnitude; otherwise the comparison
would confound horizon length with training-set size.

%---------------------------------------------------------------------------------------
\section{Hyperparameter tuning}
\label{sec:exp-tuning}

Until this point the LightGBM hyperparameters had never been tuned; they were hand-set
values. Two rounds of search were run, on the four cross-validation folds only.

\subsection{Round one: a random sample around the hand-set values}

Six parameters were varied over three levels each, centred on the hand-set values. The
full grid is $3^6 = 729$ configurations at three to five minutes per configuration per
fold, which is not affordable, so the search was staged: sixteen configurations were
sampled at random and screened on two folds chosen for maximum seasonal contrast, and the
best three plus the incumbent were then re-run on all four folds.

\begin{table}[htp]
\centering
\begin{tabular}{lll}
\toprule
{\bf Parameter} & {\bf Levels} & {\bf Controls} \\
\midrule
\verb+learning_rate+ & 0.03 / 0.06 / 0.10 & step size per tree \\
\verb+num_leaves+ & 63 / 127 / 255 & model complexity ceiling \\
\verb+min_data_in_leaf+ & 100 / 200 / 400 & minimum support per leaf \\
\verb+feature_fraction+ & 0.70 / 0.85 / 1.00 & features sampled per tree \\
\verb+bagging_fraction+ & 0.70 / 0.85 / 1.00 & rows sampled per tree \\
\verb+lambda_l2+ & 0.5 / 1.0 / 5.0 & L2 penalty \\
\bottomrule
\end{tabular}
\caption[Search space, round one]{\textbf{Round-one search space.} The middle level of
each parameter is the hand-set value, which was included in the sample so that the gain
from tuning could be read directly.}
\label{tab:space1}
\end{table}

\subsection{Round two: extending the grid outwards}

The three best configurations of round one all sat on a boundary of the grid, which means
the optimum may lie outside it. Round two extended the two directions that showed a real
effect and locked the four that did not, testing thirteen configurations.

Locking those four required an argument, because a naive reading of the round-one data
suggested otherwise. The marginal mean of \verb+bagging_fraction+ showed a monotone
0.71~\pp{} effect favouring 1.00. That effect is an artefact of the random sample: the
value 0.70 occurred only alongside the worst learning rate and the value 1.00 never
occurred alongside it, so the apparent effect is the learning rate leaking through. Within
the one homogeneous group large enough to check, the effect shrinks to 0.13~\pp{}. The
lesson --- that marginal means from a small random sample are not interpretable until
co-occurrence has been checked --- is taken up in Chapter~\ref{ch:discussion}.

\subsection{A confound that had to be removed first}

Early stopping used a fixed patience of 100 rounds. At a learning rate of 0.005 each round
advances one sixth as far as at 0.03, so a fixed patience is a systematically stricter
stopping rule at low learning rates: a configuration could look bad merely because it was
cut off early. Since the entire purpose of round two was to push the learning rate down,
the patience was scaled in proportion to $1/\eta$, anchored so that configurations at or
above the round-one winner's learning rate keep the original patience of 100 and remain
comparable with round one.

In round one the same fixed patience had penalised low learning rates, and the low
learning rate won regardless; that bias therefore made the round-one conclusion
conservative rather than optimistic.

\subsection{The early-stopping protocol mismatch}
\label{sec:exp-esmismatch}

Round two produced a configuration that won on cross-validation and then lost when
deployed. Diagnosing that led to the most consequential experiment in this group.

The tuning scripts and the delivery script early-stop on different data. The delivery
script stops on a separate validation block of 105 origins; the tuning scripts stop on the
last 15\% of the training window, roughly 970 origins. Both are defensible, but they are
not the same rule, and the discrepancy grows as the learning rate falls, because a small
stopping set gives a noisy validation curve that triggers earlier.

The quantity that makes this legible is the product of the selected number of rounds and
the learning rate, which measures how far the model actually trained:

\begin{table}[htp]
\centering
\begin{tabular}{llrr}
\toprule
{\bf Protocol} & {\bf Learning rate} & {\bf Rounds} & {\bf Rounds $\times \eta$} \\
\midrule
Cross-validation & 0.03 & 3,695 & 110.8 \\
Cross-validation & 0.005 & 25,359 & 126.8 \\
\midrule
Delivery, validation block & 0.03 & 1,530 & 45.9 \\
Delivery, validation block & 0.005 & 6,153 & 30.8 \\
\midrule
Delivery, training-window tail & 0.03 & 3,296 & 98.9 \\
Delivery, training-window tail & 0.005 & 18,556 & 92.8 \\
\bottomrule
\end{tabular}
\caption[Effective training budget]{\textbf{Effective training budget under each
early-stopping protocol.} Under cross-validation the two learning rates receive
comparable budgets, so the comparison is between learning rates. Under the original
delivery protocol the lower learning rate receives a third less budget than the higher
one, so the comparison is confounded with how far each model trained.}
\label{tab:budget}
\end{table}

An experiment was therefore added: align the delivery script's early stopping with the
cross-validation rule and re-train both configurations. Aligning the two rules has a
second effect worth noting in advance. Under the original rule the validation period was
used both to stop training and to produce the residuals from which the MinT covariance is
estimated, which makes those residuals mildly optimistic. Under the aligned rule the model
that produces the validation forecasts has never seen the validation period at all, so
those residuals are genuinely out of sample.

\subsection{Discipline on test-period use}

Across the whole project the test period was evaluated for five model versions: before
tuning, the round-one winner, the round-two winner, and those two winners again under the
aligned early-stopping rule. The first three are required by the protocol. The last two
were motivated entirely by Table~\ref{tab:budget}, which is diagnosable without reference
to the test period, and not by searching for a version that scored better. No further
test-period evaluation was performed after the fifth.
